\section{Marco metodológico}
\label{sec:metodologia}

\subsection{Metodología adoptada y su justificación}

Se adopta \textbf{PPDIOO} como ciclo de vida de ingeniería. Se trata de la metodología de ciclo de vida
de redes formulada por Cisco Systems, estructurada en seis fases: \emph{Prepare}, \emph{Plan},
\emph{Design}, \emph{Implement}, \emph{Operate} y \emph{Optimize}.

La elección se fundamenta en tres razones:

\begin{enumerate}
\item Es el marco propio del dominio de redes, que constituye el eje técnico del trabajo.
\item Corresponde al perfil de especialidad del profesor guía.
\item Sus fases coinciden naturalmente con las cuatro etapas de evaluación del curso, según muestra la
tabla~\ref{tab:ppdioo}.
\end{enumerate}

\begin{longtable}{@{}C{1.2cm}L{3.3cm}L{5.0cm}L{4.2cm}@{}}
\filaencab \encab{Etapa} & \encab{Fase PPDIOO} & \encab{Actividad principal} & \encab{Entregable} \\ \midrule \endhead
\textbf{1} & Prepare, Plan & Levantamiento del problema, requerimientos y planificación & Este documento y el backlog \\
\textbf{2} & Design & Arquitectura, modelo de datos, especificación del protocolo & Diagramas y especificación técnica \\
\textbf{3} & Implement & Construcción del prototipo y pruebas iniciales & Prototipo funcional e informe de pruebas \\
\textbf{4} & Operate, Optimize & Validación en condiciones realistas y optimización & Informe final y validación \\
\bottomrule
\caption{Correspondencia entre las fases de PPDIOO y las etapas de evaluación.}
\label{tab:ppdioo}
\end{longtable}

\subsection{Descarte de Scrum y conservación de sus artefactos de planificación}

Se descarta \textbf{Scrum como marco de gestión de equipo}. El proyecto es unipersonal, de modo que sus
ceremonias --- reunión diaria de sincronización, retrospectiva grupal, y los roles de Dueño de Producto
y Facilitador --- carecen de sentido: no existe un equipo al que coordinar.

Sin embargo, \textbf{se conservan sus artefactos de planificación}, en particular el \textbf{backlog}
con épicas, historias de usuario, tareas técnicas, criterios de aceptación y estimaciones.

\begin{cajaclave}[title={Formulación de la decisión metodológica}]
\textbf{PPDIOO define qué fases existen y en qué orden}, con el rigor propio de la ingeniería de redes.
\textbf{El backlog define cómo se prioriza y cómo se mide el avance.} Son marcos complementarios, no
alternativos: lo que se descarta de Scrum son sus ceremonias de coordinación de equipo, no sus
instrumentos de planificación.
\end{cajaclave}

La conservación del backlog no responde únicamente a una preferencia metodológica. La pauta de
evaluación asigna el 20\% de esta etapa a su construcción, y mide el avance de las etapas siguientes
como porcentaje de tareas del backlog completadas. Sin él no existiría instrumento de medición.

\subsection{Calendario y cadencia de trabajo}
\label{sec:calendario}

\subsubsection{Cadencia}

El avance se presenta al profesor guía los \textbf{miércoles}, y cada revisión se registra en bitácora
el mismo día, junto con los cambios derivados de la retroalimentación recibida. Este registro responde
a que las etapas 2, 3 y 4 asignan un 20\% cada una al criterio de incorporación de la
retroalimentación, exigiendo \enquote{ejemplos específicos de cambios en el proyecto}.

\subsubsection{Fechas de las etapas}

\begin{cajahallazgo}[title={Ambigüedad no resuelta en el calendario}]
La pauta de evaluación disponible está fechada el \textbf{13 de agosto de 2025} y sus hitos corresponden
a ese año. Existen dos lecturas posibles para el calendario 2026, que difieren en dos semanas y que en
ambos casos resultan consistentes con la regla de la pauta (último viernes de la cuarta semana del
sprint). \textbf{La confirmación con el profesor guía es la primera gestión pendiente.}
\end{cajahallazgo}

\begin{longtable}{@{}L{4.6cm}C{2.6cm}C{2.6cm}C{2.6cm}C{2.4cm}@{}}
\filaencab \encab{Escenario} & \encab{Etapa 1} & \encab{Etapa 2} & \encab{Etapa 3} & \encab{Etapa 4} \\ \midrule \endhead
\textbf{A}: la pauta se replica al calendario 2026 & \textbf{28 ago} & 25 sep & 23 oct & 20 nov \\
\textbf{B}: el sprint 1 inicia el 17 de agosto & 11 sep & 9 oct & 6 nov & 4 dic \\
\bottomrule
\caption{Escenarios de calendario. Todas las fechas caen en viernes.}
\label{tab:escenarios}
\end{longtable}

\textbf{Decisión de planificación:} se planifica contra el \textbf{escenario A}, por ser el más
exigente. Si el calendario resulta ser el B, se obtienen dos semanas de holgura; la decisión inversa
habría implicado incumplir la entrega. Un antecedente favorece al escenario A: la pauta de 2025 fijaba
la Etapa 1 el 29 de agosto, lo que implica un inicio de semestre a comienzos de agosto y no a mediados.

\subsubsection{Presupuesto de tiempo}

El estudiante desempeña una jornada laboral completa de lunes a viernes, de modo que el tiempo dedicable
al trabajo de título corresponde a tardes de días hábiles y bloques de fin de semana, estimado en
\textbf{aproximadamente 15 horas semanales}. Esta restricción condiciona el dimensionamiento del alcance
y se verifica cuantitativamente en la sección~\ref{sec:viabilidad}.

\subsection{Gestión de la configuración y trazabilidad}

Todo el trabajo se versiona en un repositorio git privado, con mensajes de commit descriptivos que
documentan no solo qué cambió sino por qué. Los documentos de trabajo se mantienen en formato de texto
plano dentro del mismo repositorio, de modo que el historial de versiones constituye evidencia directa
del avance y de la incorporación de la retroalimentación recibida.

\subsection{Riesgos identificados}

\begin{longtable}{@{}L{1.1cm}L{4.4cm}L{2.0cm}L{1.9cm}L{4.0cm}@{}}
\filaencab \encab{} & \encab{Riesgo} & \encab{Impacto} & \encab{Probab.} & \encab{Mitigación} \\ \midrule \endhead
\textbf{R1} & No disponer del hardware de sensores a tiempo & Alto: bloquea OE5 y OE6 & Media & Emulador de nodo por software; el sensor físico demuestra el caso de uso, el emulador demuestra el comportamiento bajo carga \\
\textbf{R2} & No lograr acceso a las partes interesadas & Alto: compromete el 10\% del criterio 1 y el respaldo del criterio 3 & Media & Escalar la solicitud mediante el profesor guía; documentar el intento y recurrir a fuentes secundarias \\
\textbf{R3} & Ambigüedad en las fechas de las etapas & Alto: invalida la planificación & \textbf{Alta} & Planificar contra el escenario más exigente; confirmar con el profesor guía \\
\textbf{R4} & Alcance excesivo para el tiempo disponible & Medio: riesgo de no superar el 70\% de logro & \textbf{Confirmado} & Separación entre núcleo comprometido y alcance condicional; verificación de holgura al cierre de cada etapa (sección~\ref{sec:viabilidad}) \\
\textbf{R5} & Contact ID insuficiente para telemetría continua & Medio & Identificado & Arquitectura de dos planos, documentada como decisión de diseño (sección~\ref{sec:dosplanos}) \\
\bottomrule
\caption{Registro de riesgos del proyecto.}
\label{tab:riesgos}
\end{longtable}
